% SPDX-License-Identifier: Apache-2.0
\section{Case Study: Three Failures of Self-Validation}
\label{sec:case-study}

The three episodes below have a common shape. In each case the artefact passed
every check the project had, and the check was the problem. In each case an
external reference caught it. And in each case the failure was not a mistake in
the work but a mistake in the \emph{evidence about} the work, which is the
failure mode that matters when the evidence is what you are publishing.

\subsection{Case 1: A build that succeeds because it is allowed to}
\label{sec:case-study-quick-and-dirty}

\subsubsection{The observation}

The aligned Isabelle session builds cleanly:

\begin{lstlisting}
isabelle build -D aligned
Finished VWZ_A3_Aligned (0:00:07 elapsed time)
\end{lstlisting}

Exit code 0. If this were the whole story, one would report the development as
verified.

It is not. The session's \texttt{ROOT} file carries the option
\texttt{quick\_and\_dirty}, which admits \texttt{sorry} as a proof step instead
of rejecting it. A \texttt{sorry} is Isabelle's explicit placeholder for an
unproved goal; under normal options it makes the build fail, which is exactly
what makes it useful. Under \texttt{quick\_and\_dirty} it is accepted silently.

So \texttt{Finished} means: \emph{every step that is not marked \texttt{sorry}
type-checks}. It does not mean the theorems are proved. The two statements
produce identical output.

\subsubsection{Why this is not a small caveat}

A checker that accepts everything and a correct proof are indistinguishable from
the outside. That is the defining hazard of machine-checked mathematics: the
tool's authority is exactly what makes its silence uninformative.

The consequence for reporting is concrete. \textbf{The exit code is not a
progress metric.} Any paper, README, or dashboard that reports build status as
evidence of verification is reporting a number that cannot move when proofs are
removed. Our own README now states this at the top, before the layout section,
because a reader who learns it after reading the results has already formed a
false impression.

\subsubsection{The negative control}

The remedy is a session that \emph{must} fail. We maintain
\texttt{neg/VWZ\_A3\_Neg.thy}, which states a deliberately false claim:

\begin{lstlisting}
isabelle build -D neg    # expected: exit 1
\end{lstlisting}

It fails, with exit code 1, as designed. If it ever succeeds, the session
configuration is broken and every successful build in the repository becomes
worthless.

This is the cheapest available test of whether the verification means anything,
and it is the one most likely to be omitted, because a project that is working
does not naturally generate a test that must fail.

\subsubsection{The progress metric we use instead}

Progress is counted, not inferred:

\begin{lstlisting}
node tools/count-obligations.cjs
\end{lstlisting}

The script parses the sources, counts \texttt{lemma}, \texttt{theorem} and
\texttt{sorry} per tier, and separates \texttt{sorry} inside the obligation
ledger from \texttt{sorry} elsewhere. Anything unexpected outside the ledger
fails the check.

This matters because the counter is itself a claim. An earlier version simply
failed whenever any \texttt{sorry} sat outside the ledger, which would have
forced a choice between hiding the open theorem and loosening the check for
everything. The current version names the one known open item explicitly, so
that a \emph{new} \texttt{sorry} anywhere still fails. Naming an exception keeps
the check strict; a blanket allowance does not.

\subsection{Case 2: The CBOM that failed its own standard}
\label{sec:case-study-cbom}

\subsubsection{The observation}

The project had 18 internal tests for the CBOM generator, all passing. The CBOM
was, by that measure, correct.

It was not valid CycloneDX 1.6. Validating against the official schema produced
51 errors in 7 categories: three properties that each of the twelve components
carries in the wrong place, one enumeration error on each component's
\texttt{primitive}, and three document-level keys. The pre-repair BOM (twelve
components, 102{,}895 bytes, serial number beginning
\texttt{urn:uuid:a134990a}) was generated rather than committed, so that count is
a measurement rather than something a reader can re-run from this repository; the
artefact that \emph{is} committed, \texttt{tools/cbom-cyclonedx.json} (nine
components, 63{,}485 bytes), is the repaired one and validates clean against the
vendored 1.6 schema at commit \texttt{4064ba8}.

\subsubsection{The root cause is the test, not the code}

The internal tests fed the generated CBOM back into the project's own parser and
checked that it could be read. This is circular: it verifies that the output
matches \emph{our understanding of the format}, not that it matches the format.
A parser written from the same misunderstanding will accept the same error
happily.

The errors were of three kinds, and the taxonomy is worth stating because it
generalises:

\begin{enumerate}[leftmargin=*]
  \item \textbf{Wrong location.} Risk metadata placed at
        \texttt{cryptoProperties.quantumSecurity}, a key the schema does not
        have; the correct place is a property, and in 1.6 the nearest standard
        field is \texttt{nistQuantumSecurityLevel}, which is an integer and not
        the same thing at all.
  \item \textbf{Wrong type or enumeration.} \texttt{certificationLevel} given as
        a string where the schema requires an array, and with a value
        (\texttt{known-answer-tested}) that is not in the enumeration;
        \texttt{primitive} in upper case where every legal value is lower case.
  \item \textbf{Semantic mismatch.} \texttt{implementationPlatform} set to
        ``javascript'', which is a language, not a platform. The value was true;
        the field was wrong. The fix was to move it to a namespaced property
        rather than to find a platform that happened to be close.
\end{enumerate}

\subsubsection{The general lesson}

Internal consistency is not correctness. A test suite that shares an assumption
with the code it tests cannot detect that assumption being wrong. The only
remedy is an external reference: the official schema file, the official CLI, or
a third party's fixture.

The corollary is uncomfortable for anyone shipping standard-format artefacts. If
the only validator you run is your own, you do not know whether your output is
conformant. You know only that it is self-consistent, which is a weaker and much
more common property.

\subsection{Case 3: A page that never worked}
\label{sec:case-study-viewer}

\subsubsection{The observation}

The project publishes a CBOM viewer: an HTML page that renders the bill of
materials as a sortable table with risk badges and expandable file lists. It had
been deployed, linked from the documentation, and described as working.

It had never rendered its table. The page loaded, the summary cards displayed
numbers, and the table beneath them was empty. The first call to the render
function threw:

\begin{lstlisting}
TypeError: a.toLowerCase is not a function
\end{lstlisting}

The render function iterated over a list whose entries are objects and called a
string method on each. The exception aborted the script before the table body
was populated.

\subsubsection{Why it survived}

Three reasons, and each is a general failure mode rather than a local mistake.

First, \emph{partial success looks like success}. The cards were populated by
code that ran before the throw, so the page showed real numbers. A page with
numbers on it does not look broken.

Second, \emph{the defect was downstream of the thing everyone was watching}. The
project knew the risk field had moved and expected that to be the bug. It was a
bug. But the render function never got far enough to read the risk field, so the
known defect was invisible behind an unknown one. Investigating the suspected
cause would have found a second, deeper cause.

Third, \emph{code review does not execute}. The defect is not subtle on
inspection, but inspection is not what happened; the code was reviewed for
whether the change looked right, and it looked right.

\subsubsection{How it was actually found}

By running the page's own script against a stub DOM, rather than by reading it:

\begin{lstlisting}
const scripts = [...html.matchAll(/<script[^>]*>([\s\S]*?)<\/script>/g)];
const main = scripts.find(s => s.includes('function renderTable'));
vm.runInContext(main, vm.createContext({ document: stubDocument }));
\end{lstlisting}

This takes a few seconds and immediately produced the exception above. It also
made the rest of the audit possible, because once the script runs, its output can
be compared against the source of truth:

\begin{center}
\begin{tabular}{@{}lll@{}}
\toprule
Quantity & Page claims & CBOM says \\
\midrule
Components & 12 & 9 \\
Files & 370 (hardcoded) & 304 distinct (647 occurrences) \\
Quantum-safe & 6 & 7 \\
Quantum-vulnerable & 2 & 2 \\
Certification & ``KAT passed'' for 5+ hits & \texttt{none} for all \\
\bottomrule
\end{tabular}
\end{center}

Every row was wrong except the vulnerable count, and the certification row was
not wrong so much as invented: the page derived ``KAT passed'' from an occurrence
count above five. That is a certification claim manufactured from a line count.

\subsubsection{The finding that generalises}

The page was stale in a way that no amount of care in the \emph{generator} would
have fixed, because it was hand-maintained and embedded its own copy of the data.
It displayed one document and offered a different one for download.

The structural remedy is to make drift impossible rather than to detect it: the
embedded payload should be generated from the same artefact the page claims to
display, and a check should fail when they differ. Detection by discipline works
until someone is tired.

\subsection{What the three cases have in common}
\label{sec:case-study-common}

\begin{center}
\begin{tabular}{@{}p{3.1cm}p{4.3cm}p{4.6cm}@{}}
\toprule
Case & Passed & Caught by \\
\midrule
\texttt{quick\_and\_dirty} build & the build itself & a session that must fail \\
CBOM schema validity & 18 internal tests & the official schema \\
Viewer rendering & code review, deployment & executing the page \\
\bottomrule
\end{tabular}
\end{center}

In each row, the check that passed and the check that caught the defect differ in
one respect: the passing check was internal to the project, and the catching
check was external to it. The build was the project's own build; the negative
control is an artefact whose only purpose is to be wrong. The tests used the
project's parser; the schema is the standard's. The review read the project's
code; execution ran it.

We do not think this is a coincidence, and we do not think it is specific to
this project. Self-validation cannot detect a shared false assumption, by
construction. The evidence a PQC migration needs is evidence an auditor will
accept, and an auditor is by definition external. Building the external reference
into the development, rather than adding it at the end, is the difference between
evidence and assertion.

\subsection{The \texttt{nat} subtraction obstruction}
\label{sec:case-study-nat}

One obligation deserves separate treatment, because it is a case where the
verification tool was right, the mathematics was right, and the proof was
blocked by something neither. It was open when this paper was first drafted and
has since been closed; both halves are instructive.

The obligation (Section~\ref{sec:formalization}) is
$\det(\texttt{row\_add\_matrix}\ m\ i\ j\ f) = 1$ for $i \neq j$. The proof
strategy is to expand along the first row and induct on $m$, which requires
shifting the indices down by one. The needed shift lemma is:

\begin{lstlisting}
map (%c. if Suc k = i & Suc c = j then f
         else (if Suc k = Suc c then 1 else 0)) [0..<m-1]
=
map (%c. if k = i-1 & c = j-1 then f
         else (if k = c then 1 else 0)) [0..<m-1]
\end{lstlisting}

It is false at $j = 0$. In \texttt{nat}, subtraction saturates:
$0 - 1 = 0$. So the right-hand side degenerates to
\texttt{if k = c then 1 else 0}, while the left-hand side is identically zero
because \texttt{Suc c = 0} is unsatisfiable. A proof that applies this shift
uniformly over $j$ cannot be sound, and any attempt to force it would produce
either a failed proof or, worse, a proof of a weakened statement that no longer
implies the goal.

\subsubsection{How the obstruction was found}

Before attempting the proof we verified the statement numerically, by evaluating
the determinant exactly as defined --- Laplace along the first row,
\texttt{nat} subtraction, integer entries --- over every $(i,j)$ with
$i \neq j$ for $m \in \{1,2,3,4,5\}$ and several values of $f$, all giving $\det = 1$. A control sweep with $i = j$ gives $\det = f$,
which is what it must be, and which shows the sweep is not passing vacuously.

So the mathematics was never in doubt and the missing piece was the proof. The
obligation was recorded with an explicit \texttt{sorry} and a \texttt{TODO(prove)}
note naming the false shift lemma, because the note is the useful part: the next
attempt should not rediscover the dead end.

\subsubsection{How it was closed}

The diagnosis was that the proof must split on $j = 0$ \emph{before} the shift.
That is what closed it. The final proof inducts on $m$ and splits four ways:
indices beyond the matrix in either coordinate, where the row operation does not
fire and the matrix is the identity (\texttt{aux\_t2\_nofire\_i},
\texttt{aux\_t2\_nofire\_j}); the case $i = 0$, discharged separately as the
theorem \texttt{aux\_el\_det\_elementary\_i0}; and the case $0 < i$, where the
determinant is reduced to the head-deleted minor by
\texttt{aux\_t2\_head\_reduce} and then split on $j = 0$. The $j = 0$ branch
collapses to an identity minor (\texttt{aux\_t2\_minor0\_j0}); the $j > 0$ branch
is the only place the shift is applied, and there it is valid
(\texttt{aux\_t2\_minor0\_jpos}), reducing the goal to the induction hypothesis
at $(i-1, j-1)$.

The obligation now carries no \texttt{sorry}. We report the before and the after
rather than only the after, because the value of the episode is the diagnosis:
the boundary split that the failed attempt predicted turned out to be the whole
proof.

\subsubsection{The general rule}

\textbf{Any lemma that subtracts one from an index must be split on the zero
case.} In \texttt{nat}, the intuition that subtraction undoes addition is
false at the boundary, and the failure is silent: the lemma is not rejected, it
is simply untrue, and a proof that appears to work is proving something else.

This is a small instance of the paper's theme. The type checker accepted the
statement; the statement was false; and only a counterexample at the boundary
revealed it. A formalization that never tested its own auxiliary lemmas at the
boundary would have carried the error forward into a proof that looked complete.
